\section{Implementation}
\label{sec:implementation}

The primitive lives in three Rust components: a runtime kernel that exposes a pre-dispatch admission hook, a federation crate that implements the strict bilateral DSSE verifier, and a selective-disclosure module that projects canonical receipt bytes into a BBS message vector for presentation. The wire format is the load-bearing artifact; the implementation is one realization of it.

\paragraph{Pre-dispatch admission hook.} The runtime intercepts every tool-call before dispatch through a single hook whose only output is admit-or-deny. A request without a Chiodos admission context flows through the existing non-federated path; once a context is present, three behaviours become obligatory. First, the treaty reference is resolved by reading verifier-owned state alone, so any request-borne payload purporting to declare trust roots, peer directories, signing keys, or dynamic trust bundles is treated as adversarial input rather than configuration. Second, only the in-band DSSE evidence is verified, and only against signer sets that match the treaty's declared participants under the independence requirement on public keys. Third, on any denial or aborted continuation, reserved kernel state is released back to free, so that an attacker watching a continuation get admitted then aborted cannot replay it under a fresh request.

Each of the three is a place where naive admission middleware degrades into advisory logging. Admitting federated origin without treaty context would bypass the primitive of \S\ref{sec:predicate} entirely; allowing request metadata to nominate the trust root would invalidate the verifier-owned-store assumption underwriting the freestanding accept-set theorem of \S\ref{sec:formal}; leaving continuation state allocated on denial would manifest as the replay edge the primitive is meant to close. The hook refuses each pattern as a structural matter, not as a runtime check an operator could disable.

\paragraph{Strict bilateral verifier.} \codepath{verify_chiodos_dsse_envelope} (\codepath{crates/chio-federation/src/bilateral_dsse.rs:996}) is the cryptographic primitive at the heart of this paper: a total function from a candidate envelope and verifier-owned state to either one of the six post-signature rejection codes of \S\ref{sec:predicate} or an admission verdict that binds the two named cosigners to canonical bytes. Its internal gates split into two layers. The envelope layer enforces structural properties of the DSSE bytes themselves (payload-type equality against the bilateral statement type, signature count equal to two, canonical-JSON equivalence between the inner payload and its declared digest, predicate-type equality against \codepath{PREDICATE_TYPE_CHIODOS_BILATERAL}, single-subject presence, keyid-distinctness across the two signatures). The operational layer at \codepath{crates/chio-federation/src/bilateral_verifier.rs} then accepts only envelopes that already cleared the envelope layer and gates them against state the freestanding theorem leaves as trust-store inputs: peer key pinning, ladder-manifest freshness, revocation epoch, receipt resolution, request and tool-argument hashes, policy verdict agreement, capability lease, receipt-backed-class governance receipts, and treaty binding identifiers. The ordering is load-bearing: every byte-level check fires before the verifier consults state the sender cannot see, so a sender cannot satisfy a strict subset of the gates by tailoring bytes alone.

Each of the six rejection codes traces back to a specific gate position, and Table~\ref{tab:rejection-codes} maps the codes to their source locations. The taxonomy lives on the protocol surface, not in debugging output, because an audit consumer must be able to distinguish a missing verifier-owned trust binding from a wrong predicate type from a degenerate signer reuse from a stale lease from a sibling-treaty digest disagreement; each names a different constitutional precondition with a different remediation. The threat-model scope of the verifier is also worth naming explicitly. The construction binds two honest cosigners to canonical bytes against an arbitrary network adversary; a single actor who controls both signing keys, whether through cross-party key-custody compromise or through a single vendor operating both ends of a bilateral closure, sits outside the threat model. That case is treated as an operational assumption in \S\ref{sec:limits} and is independent of every gate in the verifier.

\begin{table}[t]
\scriptsize
\caption{Six rejection codes mapped to verifier source locations.}
\label{tab:rejection-codes}
\begin{tabularx}{\columnwidth}{@{}lX@{}}
\toprule
Rejection code & Verifier source \\
\midrule
noncanonical-payload & canonical-bytes equality gate inside \codepath{verify_chiodos_dsse_envelope} (\codepath{bilateral_dsse.rs:996}) \\
predicate-type-mismatch & \codepath{predicate_type} equality against \codepath{PREDICATE_TYPE_CHIODOS_BILATERAL} inside \codepath{verify_chiodos_dsse_envelope} (\codepath{bilateral_dsse.rs:996}) \\
signer-reuse & keyid-independence gate and \codepath{require_unique_signature_keyids} inside \codepath{verify_chiodos_dsse_envelope} (\codepath{bilateral_dsse.rs:996}) \\
trust-store-miss & verifier-owned key and treaty membership lookup inside the operational verifier (\codepath{bilateral_verifier.rs}) \\
stale-lease & capability-lease freshness gate inside the operational verifier (\codepath{bilateral_verifier.rs}) \\
subject-digest-mismatch & binding-tuple subject equality gate inside the operational verifier (\codepath{bilateral_verifier.rs}) \\
\bottomrule
\end{tabularx}
\end{table}

\paragraph{Selective disclosure.} The selective-disclosure module signs a BBS projection of the canonical receipt body and verifies derived disclosure proofs against an issuer-owned key registry~\cite{boneh2004bbs,bbsDraft}. The Ed25519 signature over canonical receipt bytes remains authoritative for audit corpora; BBS is a presentation-only commitment that lets a holder reveal a subset of receipt fields to a verifier without exposing every field. The verifier pins the issuer key and the projection version, so the proof is not free-form redaction.

\paragraph{Stub disclosure.} One platform-specific stub the primitive does not depend on is worth naming. The macOS Endpoint Security path used for kernel-side telemetry is a stub: the integration surface compiles and tests pass against canned fixtures, but live Endpoint Security event delivery is not wired up. The bilateral-DSSE primitive, the admission hook, the federation crate, and the selective-disclosure module operate on canonical bytes and trust-store state regardless of whether kernel-side telemetry is live; the stub is named so a deployer planning kernel-event-driven receipts knows where the gap sits.
